\documentclass[11pt]{article}
\input{../../shared/project_style.tex}
\rhead{Basics 11 Textbook Draft}

\begin{document}
\DocTitle{Basics 11: Electric and Magnetic Fields}{II - Electromagnetism foundations}{Forces, potentials, flux, energy, and the field language used throughout plasma physics}

\begin{overviewbox}
\textbf{Textbook draft, version 1.0.} Electric and magnetic fields are the immediate agents that accelerate charged particles and transmit forces through a plasma. This chapter begins with Coulomb's law, develops electric potential and magnetic flux, and explains why field energy and field-line geometry are useful without treating field lines as material objects. This chapter is designed for a reader with no prior plasma-physics training. It distinguishes exact identities from physical approximations and connects the central equations to later simulation modules.
\end{overviewbox}

\ProjectTOC

\section{Learning objectives}
After completing this note, the reader should be able to:
\begin{itemize}[leftmargin=1.6em]
\item Distinguish force, field, potential, flux, and energy density.
\item Derive the electric field from a scalar potential in electrostatics.
\item Use the Lorentz force to interpret electric and magnetic effects.
\item Compute electric and magnetic flux through simple surfaces.
\item Check electromagnetic expressions using SI dimensions.
\item Explain the limits of field-line pictures.
\end{itemize}

\section{Prerequisites and reading strategy}
\begin{itemize}[leftmargin=1.6em]
\item Basics 1-3.
\end{itemize}
On a first reading, emphasize physical meaning and dimensional checks. On a second reading, reproduce the derivations. Attempt the computational laboratory only after the limiting cases are understood.

\section{From forces between particles to fields}

For two point charges separated by $\mathbf r=\mathbf x-\mathbf x'$, Coulomb's law is
\begin{equation}
\mathbf F_{12}=\frac{1}{4\pi\epsilon_0}\frac{q_1q_2}{r^3}\mathbf r.
\end{equation}
The electric field is defined as force per unit test charge:
\begin{equation}
\mathbf E(\mathbf x)=\lim_{q_t\rightarrow0}\frac{\mathbf F}{q_t}.
\end{equation}
The limit emphasizes that the test charge should not alter the source configuration. Superposition gives
\begin{equation}
\mathbf E(\mathbf x)=\frac{1}{4\pi\epsilon_0}\int
\rho_q(\mathbf x')\frac{\mathbf x-\mathbf x'}{|\mathbf x-\mathbf x'|^3}\,d^3x'.
\end{equation}
A field is therefore a rule assigning a vector to every point, not a stream of material moving along drawn lines.


\section{Electric potential and electrostatic energy}

In electrostatics, $\nabla\times\mathbf E=0$, so the field can be written
\begin{equation}
\mathbf E=-\nabla\phi.
\end{equation}
The potential difference is the work per unit charge required to move slowly against the field:
\begin{equation}
\phi(\mathbf b)-\phi(\mathbf a)=-\int_{\mathbf a}^{\mathbf b}\mathbf E\cdot d\boldsymbol\ell.
\end{equation}
Combining Gauss's law with $\mathbf E=-\nabla\phi$ gives Poisson's equation,
\begin{equation}
\nabla^2\phi=-\frac{\rho_q}{\epsilon_0}.
\end{equation}
The electric energy density in vacuum is
\begin{equation}
u_E=\frac{\epsilon_0E^2}{2}.
\end{equation}
Only potential differences are physically observable in classical mechanics; adding a constant to $\phi$ changes no field.


\section{Magnetic field and the Lorentz force}

The force on a particle of charge $q$ and velocity $\mathbf v$ is
\begin{equation}
\mathbf F=q\left(\mathbf E+\mathbf v\times\mathbf B\right).
\end{equation}
Because $\mathbf v\cdot(\mathbf v\times\mathbf B)=0$, a magnetic field alone does no instantaneous work:
\begin{equation}
\frac{d}{dt}\left(\frac{mv^2}{2}\right)=q\mathbf v\cdot\mathbf E.
\end{equation}
Magnetic forces nevertheless redirect particles, change momentum transport, and allow a field to confine plasma. The magnetic energy density is
\begin{equation}
u_B=\frac{B^2}{2\mu_0}.
\end{equation}
This energy density later appears as magnetic pressure, while spatial variation and curvature produce magnetic tension forces.


\section{Flux, circulation, and field-line geometry}

Electric and magnetic fluxes through a surface $S$ are
\begin{equation}
\Phi_E=\int_S\mathbf E\cdot d\mathbf A,\qquad
\Phi_B=\int_S\mathbf B\cdot d\mathbf A.
\end{equation}
Flux depends on both field strength and surface orientation. Magnetic field lines are tangent to $\mathbf B$ and may be parameterized by
\begin{equation}
\frac{d\mathbf x}{d\ell}=\frac{\mathbf B}{B}.
\end{equation}
Their visual density is conventionally used to indicate field magnitude, but this is a drawing convention. Since $\nabla\cdot\mathbf B=0$, magnetic field lines do not begin or end at isolated magnetic charges in classical electromagnetism.


\section{Electrostatic and inductive electric fields}

When magnetic fields vary in time, the electric field need not be conservative:
\begin{equation}
\nabla\times\mathbf E=-\frac{\partial\mathbf B}{\partial t}.
\end{equation}
A useful decomposition is
\begin{equation}
\mathbf E=-\nabla\phi-\frac{\partial\mathbf A}{\partial t},\qquad
\mathbf B=\nabla\times\mathbf A.
\end{equation}
The first term is potential-derived; the second is inductive. Plasma models often neglect one contribution under a stated ordering, but they should never assume $\mathbf E=-\nabla\phi$ without checking Faraday's law.


\section{Scales relevant to plasma simulation}

Typical comparisons include electric and magnetic force magnitudes,
\begin{equation}
\frac{|qE|}{|qvB|}\sim\frac{E}{vB},
\end{equation}
and field energy relative to thermal pressure,
\begin{equation}
\beta=\frac{p}{B^2/(2\mu_0)}.
\end{equation}
Such ratios identify dominant physics. In numerical work, $\nabla\cdot\mathbf B$, energy balance, and unit consistency are primary diagnostics rather than cosmetic checks.


\section{Computational laboratory}
Implement Coulomb-field and potential calculations for a pair of charges. Verify numerically that $-\nabla\phi$ matches $\mathbf E$ away from the charges. Compute electric flux through nested spheres and magnetic flux through tilted planar surfaces. Compare finite-difference divergence errors under grid refinement.

\section{Summary}
\begin{itemize}[leftmargin=1.6em]
\item Fields convert source information into local force laws.
\item Electrostatic electric fields derive from a scalar potential, while inductive fields do not.
\item Magnetic fields redirect charged particles but do not directly change their kinetic energy.
\item Flux and field-line descriptions are geometric tools whose limitations must be respected.
\item Field energies and dimensionless ratios connect electromagnetism to plasma dynamics.
\end{itemize}

\SymbolsAcronymsSection
\begin{symtable}
$\mathbf E$ & electric field & Force per unit positive charge. \\
$\mathbf B$ & magnetic flux density & Field entering the magnetic Lorentz force. \\
$\phi$ & electric potential & Potential energy per unit charge. \\
$\rho_q$ & charge density & Charge per volume. \\
$\Phi_E,\Phi_B$ & field fluxes & Surface integrals of normal field components. \\
$\epsilon_0,\mu_0$ & vacuum constants & Permittivity and permeability. \\
\end{symtable}

\section{Exercises}
\begin{enumerate}[leftmargin=1.8em]
\item Derive the electric field of a point charge from its potential.
\item Show explicitly that a magnetic force performs no work on a point particle.
\item Calculate the flux of a uniform magnetic field through a tilted circular loop.
\item Use dimensions to verify the units of $u_E$ and $u_B$.
\item Explain why field lines can move in an ideal-MHD argument even though they are not material strings.
\end{enumerate}

\section{References}
\begin{thebibliography}{99}
\bibitem{ref1} G. B. Arfken, H. J. Weber, and F. E. Harris, \emph{Mathematical Methods for Physicists}, 7th ed., Academic Press (2013).
\bibitem{ref2} G. Strang, \emph{Linear Algebra and Its Applications}, 4th ed., Brooks/Cole (2006).
\bibitem{ref3} W. A. Strauss, \emph{Partial Differential Equations: An Introduction}, 2nd ed., Wiley (2007).
\bibitem{ref4} D. J. Griffiths, \emph{Introduction to Electrodynamics}, 4th ed., Cambridge University Press (2017).
\bibitem{ref5} J. D. Jackson, \emph{Classical Electrodynamics}, 3rd ed., Wiley (1998).
\bibitem{ref6} F. F. Chen, \emph{Introduction to Plasma Physics and Controlled Fusion}, 3rd ed., Springer (2016).
\bibitem{ref7} R. J. Goldston and P. H. Rutherford, \emph{Introduction to Plasma Physics}, CRC Press (1995).
\bibitem{ref8} P. M. Bellan, \emph{Fundamentals of Plasma Physics}, Cambridge University Press (2006).
\end{thebibliography}

\end{document}
